% Section 13 of "Communication, Coherence and Collective Agents".
\section{Discussion}\label{sec:dis}
% src: M_spec Section 6 (dis) and decisions 3, 4, 6, 7, 11, 12, 14, 15, 16; E_edits_digest (final statuses of all sections); H_brief Section 0 (the author's two messages); M_spec_H5 Sections 0, 2 to 4 (H5, decisions 19 to 22, Friston operator attribution); sections sec:met, sec:hie, sec:ren, sec:env, sec:rel-renormalization in final form; E3 proof checks

We return to the author's five ideas and say for each what is proved and in which models, and what is only computed, conjectured or open (\cref{sec:dis-status}); we take up the question behind the fifth idea, stating when a meta-agent is an agent with external states of its own and what a sensor optimal for a declared predictive objective reads at each level (\cref{sec:dis-meta}), and compare the coarse-graining of \cref{sec:ren} with the renormalization of Friston and colleagues (\cref{sec:dis-renormalization}); we then collect the main open problems and limitations (\cref{sec:dis-open,sec:dis-limits}) and relate the framework to the author's gauge-theoretic program (\cref{sec:dis-gauge}). A result cited with the word Supplement is a secondary statement placed in the supplement, whose numbers carry the prefix S; the proofs, the numerical checks and the extensions that address open problems posed in the literature are there as well.

\subsection{Status of the hypotheses}\label{sec:dis-status}

\Cref{hyp:int-h1} carries the first idea, that meta-agents form locally and have Markov blankets assembled from their members' active, sensory and internal states. Its boundary clause is a theorem that selects nothing. With roles assigned to channels and noise specific to each channel, every set of agents that hears its sources through separate channels is an exact path blanket from every conditionally independent start (\cref{thm:col-path-blanket}), and, when the world-conditional covariance decays exponentially and is bounded with bounded inverse on a graph of polynomial growth, every region with a finite interior whose shell is thick on the scale of an explicit bound on the Markov length is an approximate stationary blanket of the world-conditional law (\cref{thm:col-boundary}; Supplement, \cref{thm:col-thick-shell,thm:col-markov-bound}). The correlation length does not bound the Markov length (\cref{prop:col-markov-vs-correlation}), so an earlier formulation of H1, that a shell thicker than $\xi$ makes a blanket, was false (Supplement, \cref{rem:col-shell-error}). Since every candidate with separate channels has such a blanket (Supplement, \cref{prop:col-selects-nothing}), collectives are selected by contrast on the world-conditional precision at a level declared above the ratio of the largest to the mean leak on the region and above the baseline of a comparison medium that keeps the population's declared geometry (\cref{def:col-baseline}), a definition that includes the moment bound and needs a diffusive precision (\cref{def:col-collective}); diffusivity is proved for reversible dynamics with nonnegative weights and row sums below one and for nonnegative normal dynamics with isotropic noise under a condition on row and column sums, and it fails in general (\cref{thm:col-diffusive-precision}; Supplement, \cref{ex:col-non-diffusive}). Homogeneous lattices have bounded contrast, so the baseline is finite (\cref{thm:col-lattice-contrast}), and a homogeneous medium, being its own comparison medium, has no collective at any admissible level (\cref{def:col-baseline,def:col-collective}); contrast bounds extent for couplings of any range (\cref{thm:col-contrast-size}), and spectral and correlation criteria are bracketed explicitly only for two equal blocks that form the population (\cref{thm:col-bridge}). Structure that a declared geometry already contains is assumed rather than found, and a finite closed population with uniform leak is a collective of its own comparison medium only if it is more connected than that medium (Supplement, \cref{rem:col-declared-geometry,prop:col-whole-population}).

The closure, interaction and agency clauses of \cref{hyp:int-h1}(c) are theorems in narrower settings (\cref{thm:agy-h1}), on exact closure of a linear coarse variable (\cref{thm:agy-mori-zwanzig,thm:agy-closure-identity}), on the order parameter under causally separating nulls (\cref{prop:agy-nulls-vanish}) and on deliberation under an asynchronous protocol with shared data and intentions (\cref{thm:agy-agency}); a meta-agent is a definition (\cref{def:agy-meta-agent}), no population analyzed in the manuscript is shown to meet every clause, and whether the clauses can hold jointly is an open problem (Supplement, \cref{open:agy-meta-agent-exists}, with the scope of each clause in \cref{sup:dis-other-clauses-of-h1}).

\Cref{hyp:int-h2} carries the second idea, that communication is an action and receiving a sensory process. The author's maps have exact forms (\cref{prop:fw-prior-completion}): the sender's $T$ is its effective code $H_j=T_cU_j$, which keeps its mean only modulo $\ker H_j$; the receiver's $O$ is the right inverse of its code model that its own prior selects; and a correctly read noiseless message acts through a projector in the receiver's metric, so understanding is prior-dependent completion, with an exact formula for the meaning lost. Direct, stigmergic and performative channels are modes of one transfer kernel (\cref{prop:fw-transfer-kernel}). Sharpened as misreading of joint structure, H2 holds in stated scopes. Literal receivers overvalue redundant messages, which is double counting, in the scalar track model within a band of positive error covariances and in static common-source worlds that are scalar, aligned or read through invertible noiseless codes (\cref{prop:mis-overvaluation-scope}; Supplement, \cref{prop:mis-track-band}); they undervalue messages with complementary errors in the track model, and other messages in closed-form examples (Supplement, \cref{prop:mis-track-band,ex:mis-cancellation,ex:mis-sender-completion}), and elsewhere undervaluation is found by search in two random ensembles with one seed, with counts specific to those ensembles and no frequency independent of them (Supplement, \cref{obs:mis-random-codes}). The level-1 likelihood is true only when the coded error is independent of the world and of what the receiver knows, which in static worlds of exact-posterior agents a proper common prior rules out for every nonzero code when the receiver reads with the sender's code and header (\cref{prop:fw-level-optimality}). Agreement cannot reveal a shared error (\cref{lem:mis-social-blind}), closed groups of scalar linear receivers with fixed gains, homogeneous grounding and reversible weights are echo chambers (\cref{thm:mis-echo-chambers}), and exact inference forms no weak or self-sealed one (\cref{prop:mis-exact-no-echo}). Read this way, H2 makes exact fusion a clause that a collective must survive (\cref{def:agy-nulls}). Exact conditions for each direction of misreading are open (Supplement, \cref{open:mis-scope}).

\Cref{hyp:int-h3} carries the third idea, that coherent message passing is part of how meta-agents form. Coherence splits into fidelity, calibration and consistency, each scored with accuracy, because a closed group can be coherent and still be an echo chamber (\cref{def:coh-components}). On the infinite ring of scalar level-0 receivers with exact code models, information about a distant source decays at twice the mass (Supplement, \cref{prop:coh-info-decay}). Receivers with fixed assigned noise cannot withdraw from mistranslated senders. In a static two-group model at infinite policy precision, calibrated receivers with at least $\kappa$ in-group senders available withdraw from every cross-group sender whose mistranslation cost exceeds its information advantage, and the reverse inequality for a cross-group sender guarantees only that some cross-group sender stays attended (\cref{prop:coh-c0-blind,thm:coh-withdrawal}). Consensus exists exactly when its content is fixed by the holonomies of the fundamental cycles (\cref{thm:coh-holonomy}), and true translations between invertible encodings of one world are flat (Supplement, \cref{prop:coh-flat-true}), so the inconsistency that drives withdrawal is there learning error, which grounded learning with full-rank codes and convergent translations confines in the limit to edges attended finitely often (Supplement, \cref{prop:frm-flattening}). Amplification around cycles needs miscalibrated reading of senders' uncertainty (\cref{thm:coh-amplification}), and the grounding tension makes every agent's stationary mean-square error at least $Q\xi/2$ in scalar level-0 networks with exact code models, time-invariant weights and homogeneous grounding, under conditions (Lin), (Loc), (B) and (NN) of \cref{asm:sca-standing} (\cref{thm:coh-grounding-tension}). The claim that boundaries are where inconsistency and low fidelity concentrate holds in a qualified form, since withdrawal can fail across complementary access to the world (Supplement, \cref{ex:coh-unique-information}) and can miss inconsistency spread thinly along attended cycles (Supplement, \cref{prop:coh-attended-frustration}).

The formation half of H3 is a theorem only in part and only in one model. In the two-group model, two groups whose private frames differ by a rotation, with idealized calibration and attention capacity below the group size, clauses (T1) to (T3) and (T5) of the formation target (\cref{def:frm-target}) are theorems at $\gamma=\infty$ with exact in-group translations ($\sigma_\alpha^2=0$), the contrast clause at the levels of the window of \cref{thm:frm-split-contrast}(c), and there the declared start is itself the split (\cref{thm:frm-formation}(e)). For $\gamma_S<\gamma<\infty$ the theorems are reachability of the split branch and, above $\gamma_T$, a single fixed point on that branch that attracts every state, and they are theorems about the frozen-calibration mean-field map (\cref{thm:frm-formation}(a),(b)); the contrast of the reached state is computed at the base point (Supplement, \cref{obs:frm-exploration}), and the stochastic-approximation step from the full model to the map is open (\cref{open:frm-closure}), although from the declared start the full model with idealized calibration ends on the split branch in every replicate computed at $\gamma=80$ and $160$ (\cref{obs:frm-stochastic}). Formation there is a calibration-level effect of misreading joint structure. Receivers that know their translation errors pay no mistranslation cost, so from a start law exchangeable given the world path the split fails clause (T4) for the translation-aware null of type S, the exact-fusion persistence clause, which is proved (\cref{prop:frm-exact-fusion}), and the model yields no formation result in the sense of \cref{def:frm-target}, while for full cross-covariance-aware exact fusion the failure is heuristic; above the threshold $\gamma_T$ the split coincides with the partition by codes. The split differs from the deletion nulls by proof, and from the frozen-environment null only when the initial states that frozen messages carry are deterministic or conditioned on (\cref{prop:frm-nulls}(a)). Persistence under exact fusion with drifting frames is a conjecture (Supplement, \cref{conj:frm-drifting-frames}). Online calibration from a symmetric start reaches the split's attention pattern (not necessarily a split selected by contrast) in 0 to 7 of 8 replicates per cell, with reruns of a cell differing by up to three replicates, and is open (\cref{obs:frm-online}; Supplement, \cref{open:frm-online}). Formation beyond the two-group model is a conjecture (\cref{conj:frm-formation}). Coherence thus enters formation through calibration, as the author proposed, but the collectives it forms in the one model analyzed are attributed to misreading of joint structure, not to communication as action, a verdict proved for the translation-aware null and heuristic for full exact fusion.

\Cref{hyp:int-h4} carries the fourth idea, that the informational correlation length and time grow with the scale of collectives, as the telephone and international organizations suggest. In final form it concerns linear, local, bounded networks, conditions (Lin), (Loc) and (B), and, apart from the exact Ornstein--Zernike and second-moment lengths and the exact secular root and memory criterion of the record model, it is a set of upper bounds, each conditional on the moment bound (M) and a named contraction (\cref{thm:sca-h4}). Without (M), which the agents' own attention can violate (Supplement, \cref{ex:sca-amplifying-profile}), only finite-time statements such as the light cone hold (\cref{thm:sca-lightcone}). Under (M), nearest-neighbor passing on $\Z^D$ has the Ornstein--Zernike length $1/\arccosh(1+\varepsilon/(2a))$ exactly, so lengths grow like $\ell(a/\varepsilon)^{1/2}$ (\cref{thm:sca-oz}); spatial reach grows like the square root of memory (Supplement, \cref{prop:sca-diffusive-scaling}); records lengthen memory exactly under the criterion of \cref{prop:sca-records-memory}; and, with uniform leak, the largest extent obeys $\Lstar\le\xi_2/\sqrt{r_0-1}$ (\cref{cor:col-lstar}). The author's examples have exact counterparts: links of every length, which the telephone adds, make correlations decay algebraically when their weights have power-law tails (\cref{prop:sca-longrange}), with no strict light cone (Supplement, \cref{prop:sca-fronts}), and archives lengthen memory when they decay more slowly than members forget. For the number of agents coherent with one variable, the bound $L\lesssim\xi$ is right in low dimension and loose in high dimension (\cref{cor:sca-size}). None of these bounds says that collectives of that size form. Whether attention selected by expected free energy drives $a$ and $\varepsilon$ to the regime in which $\Lstar$ is comparable to $\xi_2$ is open (Supplement, \cref{open:sca-endogenous}), and the formation model has no geometry, so the third and fourth ideas do not yet meet (Supplement, \cref{open:frm-access}). The pointwise Ornstein--Zernike form in two and more dimensions rests on a restated theorem and a heuristic step (Supplement, \cref{rem:sca-oz-form}), and condensation beyond trees, the decline of shared rank with scale and the absence of amplification with one code are conjectures (Supplement, \cref{conj:sca-condensation,conj:sca-rank-down,conj:sca-common-code}).

\Cref{hyp:int-h5} carries the fifth idea, that a meta-agent is a bona fide agent with internal and external states of its own, whose external states are a sort of abstract physics, and that levels are linked by top-down priors, bottom-up evidence and renormalization. Its theorems are proved in linear-Gaussian models, apart from one extension of the code game and the ligand and sedimentation statistics of a cell (Supplement, \cref{prop:hie-record-bistability,prop:env-concentration,ex:env-cell}), for coarse states, partitions and comparison geometries declared before the analysis (\cref{asm:int-a12}). Part (a) is an exact conditional theorem (\cref{thm:met-h5a}). When, in a stable joint process with per-channel noise, the members' sensory channels are memoryless readers of the coarse exterior, summarized by the information-form statistic, when the exterior reads the members only through their summed act, when the coarse exterior has an exact abstract physics whose drift is detectable through the information matrix of that statistic, as when the drift is stable, and when the members pool information contributions by sum-preserving consensus, the coarse sensory and active states are an exact path blanket from every conditionally independent start, the coarse internal state is exactly closed given that blanket, the members' average estimate is, up to the response to their averaged computation noise, the predictive mean of the steady-state coarse filter at every number of consensus rounds, which with the steady-state covariance parametrizes the meta-agent's predictive law of its current exterior from the centralized start with stationary initial covariance, and, under the deliberation protocol of the group model with a shared preference precision, joint acts follow a group posterior with habit tilts bounded by the members' disagreement. The theorem has an approximate form for stable linear-Gaussian collectives whose members run the consensus filter: with the blanket conditions exact and the abstract physics replaced by a declared model with closure and input defects, the blanket and the closure of the internal state remain exact, and the path law of the coarse exterior, the belief and the law of the joint act carry error terms in the defects and the computation noise, with no smallness condition on the defects (Supplement, \cref{thm:apx-h5a}). Error bounds also exist for components: the leak rate of an approximate coarse blanket is bounded, for small violations of its conditions, by a quantity quadratic in the violation, with a bound that exceeds the leak by four orders of magnitude in the checked system and a constant about $20$ times the exact one (\cref{prop:met-leak-bound}; Supplement, \cref{open:met-leak-constant}); the memory of one coordinate of an approximately closed exterior with its input held fixed is bounded, in continuous time (for non-symmetric drift under two smallness conditions on the relative dispersion) and for reversible exteriors with spectrum in $[0,1)$ in discrete time, at fixed relative dispersion, by a quantity that falls like the inverse cube of its drift-based contrast (\cref{prop:met-memory}); and, under property (iv) of the theorem and with the means of the members' beliefs equal to their filter states, the bound on the habit tilts vanishes as consensus rounds grow, up to the response to the members' centered computation noises, under symmetric, doubly stochastic and nonnegative consensus with $\rho_\perp<1$ (\cref{thm:met-h5a}(v)). For the predictive law under approximate closure the stationary bound needs a stable linear-Gaussian joint process and a positive definite record noise, needs no smallness condition, includes the members' computation noise, and holds in the upper limit from every Gaussian start, also for an exterior that is stable only under the collective's acts (Supplement, \cref{thm:apx-unconditional,prop:apx-start,thm:apx-h5a}); at a finite round there is a finite-horizon bound that accumulates over the horizon (Supplement, \cref{prop:apx-posterior-finite,open:met-approximate}). For the closure of the coarse internal state under level-0 protocols there is none beyond the exact recursion of the members' average for scalar receivers (Supplement, \cref{prop:met-level0}). Because the five properties hold at every number of consensus rounds and none uses contrast, H5a describes the structure of a meta-agent and selects none (\cref{rem:met-no-selection}); which sets of agents are meta-agents is still decided by contrast at a declared level and by the null family (\cref{thm:agy-h1}).

Part (b) holds in linear-Gaussian hierarchies: on Gaussian trees exact inference sends likelihoods up and leave-one-out priors down (\cref{thm:hie-messages}), linear conformity to a self-constituted prior selects no partition that its declared couplings and readership do not contain (\cref{prop:hie-conformity}), and formation by conformity is a conjecture (Supplement, \cref{conj:hie-conformity-formation}). Part (c) is proved for linear sensors of linear-Gaussian exteriors (\cref{thm:env-canonical}), and beyond such exteriors the selection is a conjecture (Supplement, \cref{conj:env-nonlinear}). Part (d) is proved for Gaussian coarse-graining under the moment bound and a diffusive precision with positive leaks (\cref{asm:ren-standing}), and its depth bound for decimated media in two and more dimensions only under an inequality proved in one dimension and checked for one decimation step in two (\cref{prop:ren-depth}), while the decay of decimated kernels in two dimensions and the first-order rate of the memoryless reduction are numerical observations and which regions of high ratio Laplacian renormalization returns is an open problem (Supplement, \cref{obs:ren-decay,obs:ren-adiabatic-rate,conj:ren-lrg}). The statements behind the three parts are collected with their scopes in the supplement (\cref{sup:dis-parts-b-to-d-of-h5}).

\subsection{Meta-agents, their abstract physics and what each level senses}\label{sec:dis-meta}

The author asked whether a meta-agent is a bona fide agent with internal and external states of its own, its external states a sort of abstract physics, and what agents at different levels sense, from a corporation and a person to a nerve cell. In linear-Gaussian models the answers are as follows, each with the scope of the results it rests on.

A set of agents is an agent in the sense of the framework (\cref{sec:met}), with coarse internal, sensory, active and external states, when its declared coarse states meet the hypotheses of \cref{thm:met-h5a}: the members' sensory channels are memoryless readers of a few coordinates of the exterior and are summarized by the information-form statistic, the exterior sees the members only through their summed act and no sensor reads an individual member's act, the coarse exterior has an exact and detectable abstract physics, the members pool information contributions, and, for the joint act, they deliberate under the protocol of the group model with a shared preference precision. It then has an exact blanket, a closed internal state, a model of its abstract environment and a joint act that follows a group posterior, and an approximate form of the theorem is proved for stable linear-Gaussian collectives with exact blanket conditions and is open beyond them (Supplement, \cref{thm:apx-h5a,open:met-approximate}). These conditions are sufficient, and dropping any one makes a property fail by a computed amount in the instances of \cref{tab:met-hypotheses}: an equal-weight summary of the members' sensors leaks, sensors with unequal memories make the information-form statistic a collider, a product pool of the members' local filters is overconfident, at a relative entropy of $0.127$ nats, and literal pooling of posteriors from a common prior is overconfident by proof, and a firm that models the price alone when its competitor's states are needed is misspecified by $0.074$ nats (Supplement, \cref{ex:met-summary-leak,ex:met-channel-memory,prop:met-literal-aggregation,ex:met-two-firms}). Having this structure is not being a collective: it holds for sets of every contrast, and selection stays with contrast at a declared level and with the null family (\cref{rem:met-no-selection}).

The abstract physics of a meta-agent is the law of its declared coarse exterior given its coarse act. It is exact if and only if the kernel of the declared map is invariant under the exterior dynamics and the coarse exterior reads the members' channels only through the coarse act; for block sums the kernel condition is lumpability (\cref{thm:met-abstract-physics,def:met-lumpable}). A declared exterior that fails the kernel condition has a unique smallest completion that meets it, made of itself and the part of the hidden exterior that it eventually feels, of minimal order $d^\ast$ when the completion reads the members' channels only through the coarse act; uncompleted, it carries memory bounded by its coupling to the hidden exterior and a timescale ratio (\cref{thm:met-abstract-physics}(f), \cref{prop:met-memory}). Whether a meta-agent has an abstract physics is therefore a property of the declared variables, decided by closure, not a choice of description. The variables that close can include other organizations: in the two-firm model the abstract physics of one firm is a five-dimensional law of the price and of its competitor's sensory statistic, estimate, stored act and order flow (Supplement, \cref{ex:met-two-firms}). Aggregates of many agents close when the agents form interchangeable types, with stationary deviations from their infinite-population limit whose variance is at most of order one over the size of the smallest type, and keep memory when their persistences differ, however many agents there are, with long memory in the infinite-population limit under a common shock when the persistences are dense enough at the unit root (Supplement, \cref{prop:env-aggregate-closure,prop:env-aggregate-memory,ex:met-traders}); in a worked example an exterior made of another collective is more nearly closed the larger that collective's drift-based contrast (Supplement, \cref{ex:met-other-collectives}).

The author also asked whether meta-agents link levels as the hierarchical models of organisms do, a person receiving priors from the meta-agents a level above, a culture, a job or a community. In linear-Gaussian hierarchies they can, with one difference (\cref{sec:hie}). A member can carry its collective as a latent cause whose prediction is an empirical prior on its own states, and this two-level model is exact if and only if the collective coordinate is closed (Supplement, \cref{prop:hie-closure-exact}); where contrast is high and the drift reversible, a member that models the collective average as a Markov latent cause is misspecified only to fourth order in the dispersion of leakage (Supplement, \cref{prop:hie-reversible-sandwich}). On the Gaussian tree of a collective and its members, exact inference passes likelihoods, equivalently precision-weighted prediction errors, up and leave-one-out priors down, and predictive coding reaches the same means (\cref{thm:hie-messages}; Supplement, \cref{prop:hie-predictive-coding}). The difference is that the levels of a hierarchical model of an organism are parts of one generative model, while the prior that a member receives from the level above is built by other agents, often from the members' own reports, and reaches the member only through its sensory channels, as messages, records, prices and acts. The exact top-down message is therefore the collective's estimate with the member's own contribution removed; a person who reads a community's summary that includes the person's own report as independent evidence double counts, and a summary of a static quantity rebuilt every round from such readings grows overconfident without bound under the union, pooling and product rules (\cref{prop:hie-literal-once}; Supplement, \cref{prop:hie-literal-records}). A record that members write, such as an archive, a ledger or a newsletter, carries the collective's prior with a memory longer by a factor of about $1+\beta/\omega$ when it does not decay and, for grounded variables, a shared error that the members cannot detect among themselves (\cref{prop:hie-record-gain,thm:hie-floor,lem:hie-record-blind}), and linear conformity to such a self-constituted prior selects no partition that the declared couplings and readership do not already contain (\cref{prop:hie-conformity}). Reading culture, occupation and community in these terms is an interpretation (Supplement, \cref{rem:hie-gauge}).

What a sensor optimal for a declared horizon-$h$ attention objective reads is set by coupling, horizon and preferences (\cref{sec:env,def:env-attention-objective}; Supplement, \cref{prop:env-selection}); the objective, predictive information at the horizon plus preferences, is declared, and the expected free energy of the framework does not imply it and can oppose it (Supplement, \cref{rem:env-efe-horizon-free,rem:env-two-step}). A capacity-limited linear sensor optimal for predictive information reads the leading canonical variates of its exterior at its horizon, in reversible exteriors with distinct mode moduli under a budget fewer and slower ones as the horizon grows, and spatially coarse modes in diffusive exteriors, and preferences select within this order and can reverse it (\cref{thm:env-canonical,prop:env-horizon,thm:env-control}; Supplement, \cref{prop:env-diffusive}). If a meta-agent's sensory statistic is linear in its members' readings, its exterior channels fix its budget and its closure timescale sets its horizon, the same theorems say which coarse variables a sensor optimal for that objective reads, slower and coarser ones for slower meta-agents, which follows from the declared objective and horizon and is no independent law (Supplement, \cref{sec:env-levels}). The author's three examples then read as follows, each a worked model whose reading is an interpretation. A corporation senses an informational environment because at its decision horizon the predictive and relevant variables are slow aggregates: a firm that sells at $t+h$ what it produces at $t$ loses nothing by reading the futures price of maturity $h$, which is the separation theorem of production under price uncertainty in its risk-neutral form \citep{feder1980futures} and holds here by definition of the futures functional, and the only model-specific content is that the spot price serves it better than a slow index only below a crossover lag, $2.19$ rounds at the computed parameters of the declared two-mode market, with the price taken as exogenous, without price impact, and read without noise (Supplement, \cref{ex:env-firm,rem:env-firm-meta-agent,rem:met-corporation}). A person senses the economy through one or two numbers: with three decisions at horizons $1$, $8$ and $40$ and a Shannon budget, the optimal signal reads one slow aggregate up to $0.72$ nats per period and adds a channel mainly on local prices above it, in a reversible exterior noiseless aggregates of the slow modes suffice for a long-horizon decision up to an error that decays geometrically in the horizon, and physical channels can add much only at short horizons (Supplement, \cref{ex:env-person,prop:env-multiplex}); whether the economy's aggregates are closed at a person's timescale is not established. A nerve cell senses the chemical potential of its ligand and not gravity because the ligand couples to its blanket above thermal noise, is predictable at its horizon and enters its preferences, while a molecular gravity sensor reads with a signal-to-noise ratio of $2.2\times10^{-14}$ per particle and reading, against a chemical estimate that integrates about $7.9\times10^3$ molecule counts; at equal pooling the gravity ratio is $2\times10^{-10}$, nearly ten orders of magnitude below one, and gravity is absent from the cell's preferred outcomes by construction of the model (Supplement, \cref{ex:env-cell}). The gravity figure is the barometric exponent of equilibrium statistical mechanics, a known physical bound that the framework classifies and does not derive. By textbook accounts that we did not check against primary sources, cells that do sense gravity use dense bodies of micrometer size, so the statement holds for typical neurons and fails for vestibular hair cells, and an organism, a meta-agent of cells, senses gravity through such specialized members (Supplement, \cref{rem:env-neurons}).

\subsection{Relation to renormalization in active inference}\label{sec:dis-renormalization}

Friston and colleagues build scales from particular partitions with an adiabatic reduction \citep{friston2019free} and, in discrete state spaces, from renormalizing generative models \citep{friston2025pixels}. In the linear-Gaussian framework the corresponding steps have exact forms. Decimation and block averaging are exact maps of Gaussian laws, so what can fail under coarse-graining is closure: of a finite description, of the diffusive class on which contrast is defined, or of Markov dynamics (\cref{sec:ren}). The linear reversible analogue of their reduction operator keeps the slowest mode of each particle; it is exact if and only if the slow mode of each particle drives the other particles only along their slow modes, the ensemble assumption in linear form, and otherwise errs by one matrix, the static error of the Galerkin description, which equals the time integral of the exact memory kernel and is bounded by spectral quantities of the partition, by $\Vert\Pi_{\perp\Phi}\Vert^2/\xin^{\min}$, with $\Pi_{\perp\Phi}$ the coupling of the block averages to the coordinates within blocks, at most the cut dispersion $\sigma^2_{\mathrm{cut}}$ over the smallest internal rate, when the retained variables are the block averages of collectives with uniform leak (\cref{prop:ren-perron,lem:ren-galerkin,thm:ren-memory}). Their grouping operator eliminates internal states by conditioning on blankets. In linear dynamics this keeps the first term of an exact Mori--Zwanzig equation and replaces the rest by white noise, so it drops a memory kernel, an initial-state term and the color of the force, and the retained process is Markov with a transition law independent of the initial law only if the retained states do not read the internal ones, which fails for every particle whose active states read its internal states (Supplement, \cref{prop:ren-adiabatic}). The size of what is dropped is observed, not proved: in two computed instances the error falls like the inverse of the factor by which the internal states are faster, once that factor is ten or more (Supplement, \cref{obs:ren-adiabatic-rate}). In the pages of these works that we read we found no error bound for the reduction, a search result rather than a survey (Supplement, \cref{sec:rel-renormalization,rem:ren-rgm}). As we read them, renormalizing generative models build conditional independence into a discrete generative architecture rather than derive it from a given dynamics, and bounds for discrete and nonlinear coarse-graining, which would apply to them, are open here (Supplement, \cref{open:ren-beyond,open:met-nonlinear}). Two consequences of the exact forms bear on how meta-agents are judged: the exact law of block averages can leave the diffusive class, so contrast between meta-agents is computed in coarse media, and in a homogeneous medium, without contrast, the coarse media overstate the coupling between blocks by the block factor (Supplement, \cref{ex:ren-block-average,ex:ren-no-contrast}).

\subsection{Open problems}\label{sec:dis-open}

The open problems that bear most on the five ideas fall into five groups; all but \cref{open:frm-closure} are stated in the supplement. For formation, persistence under exact fusion with drifting frames (\cref{conj:frm-drifting-frames}) decides whether formation here is more than misreading; the stochastic step, online calibration and self-consistent calibration (\cref{open:frm-closure,open:frm-online,open:frm-self-consistent}) decide whether agents reach what the mean field selects; and access and geometry, with endogenous couplings and leaks (\cref{open:frm-access,open:sca-endogenous}), would join formation to scale. For the definition of a collective, contrast for non-diffusive precisions (\cref{open:col-non-diffusive}), the world-conditional precision of receivers that fuse exactly (\cref{open:col-exact-fusion}) and brackets beyond two equal blocks (\cref{open:col-bridge}) are open. For the other clauses of H1, closure in discrete time beyond the reversible case of one coordinate (\cref{open:agy-discrete-closure}), generic positivity of the order parameter under selected attention (\cref{open:agy-attention-positivity}), an agency bound uniform in policy precision (\cref{open:agy-uniform-bound}) and agency of abstract collectives (\cref{open:abs-institutions}) are open. For tests, an estimator of the order parameter without a recorded shared drive is needed before the interaction test applies to populations that act on their world: a partial route covers a static or stationary world without performative acts under sparsity and spread conditions not verified here, and in a linear-Gaussian network with an unrecorded driver positivity of the order parameter is certified when a channel intervention changes the recorded law, under the invariance conditions (I1) to (I3) of the intervention theorem, and is not settled, without restrictions on the driver, by records under one regime; the estimator itself remains open (Supplement, \cref{open:agy-unrecorded-drive,thm:idn-intervention}). For the fifth idea, an error bound for the meta-agent's predictive law under approximate closure beyond the stable linear-Gaussian case, together with a bound on the memory of the coarse internal state, would give H5a an approximate form (Supplement, \cref{open:met-approximate,thm:apx-unconditional}); the population law of the two-group split under full exact fusion would decide, beyond one receiver's deviation, whether the split survives the exact-fusion null (\cref{open:hie-exact-fusion-population}); a contrast and a comparison medium for populations with record nodes would decide whether member-written records make collectives (\cref{open:hie-record-contrast,open:col-records-baseline}); contrast for the non-diffusive laws that block averaging produces, and closure in discrete time for non-reversible updates, would let hierarchies be judged on the exact coarse law and extend the closure lemmas (\cref{open:ren-beyond}); formation by conformity would give formation without misreading (\cref{conj:hie-conformity-formation}); error bounds for discrete and nonlinear coarse-graining would reach the setting of renormalizing generative models (\cref{open:ren-beyond,open:met-nonlinear}); and the sensor design of an agent that filters its readings would extend the theory of selective sensing to agents with memory (\cref{open:env-filtering}).

The ten open problems that the recent literature on collective active inference poses are set against these results in \cref{sec:opn}, where \cref{tab:opn-status,tab:opn-status-b} list for each what is established, where, and what stays open. The extensions developed for them in the supplement narrow problems of the present list and close none: \cref{sec:apx} gives the stationary bound of \cref{thm:apx-unconditional}, the approximate meta-agent theorem built on it (\cref{thm:apx-h5a}), and a partial result on the subject of \cref{open:ren-beyond} that settles none of its items, \cref{sec:idn} treats whether positivity of the order parameter can be certified, by records and by interventions, and leaves its estimator open, and \cref{sec:prv} narrows the related problem of the provenance that a record must carry in a dynamic hierarchy (Supplement, \cref{open:hie-provenance}) by giving storage under which fusion is exact for one shared linear-Gaussian state, while the tree with a different variable at every node stays open.

\subsection{Limitations}\label{sec:dis-limits}

The theory is linear and Gaussian: lengths, contrast and closure concern second moments of linear error recursions, and the results that hold more generally do not supply the contrast that selects collectives; among them the formation theorem bears on which groups form, and it is a statement about a nonlinear mean-field map (\cref{sec:pre-scope,thm:frm-formation}). The sections on the fifth idea (\cref{sec:met,sec:hie,sec:ren,sec:env}) leave this setting only in one extension of the code game and in the ligand and sedimentation statistics of the cell (Supplement, \cref{prop:hie-record-bistability,prop:env-concentration,ex:env-cell}); discrete state spaces, nonlinear dynamics and renormalizing generative models otherwise appear only in remarks, open problems and conjectures (Supplement, \cref{open:met-nonlinear,open:ren-beyond,conj:env-nonlinear,conj:hie-conformity-formation}), and their firms and people are worked examples within it. The dynamic receivers of the theory read at level 0 or level 1. Exact fusion enters only as a reference and as a null; its world-conditional precision is not known to be diffusive (Supplement, \cref{open:col-exact-fusion}), and in the formation model only its translation-aware part is analyzed (\cref{prop:frm-exact-fusion}). Every length, time and volume assumes stationarity and the moment bound, and most theorems take attention fixed or exogenously scheduled, while the agents' attention depends on their states and can violate the bound (Supplement, \cref{ex:sca-amplifying-profile}); for state-dependent schedules the uniform bound is proved only on finite populations (\cref{lem:sca-relations}). Contrast needs a diffusive precision, which non-reciprocal attention and noisy records need not produce (Supplement, \cref{ex:col-non-diffusive,rem:col-which-contrast}). Formation is a mean-field result for one model, with frozen calibration, identical agents, two groups, rotation frames in the plane, a full-rank public code, scarce capacity and no geometry (\cref{def:frm-two-group,asm:frm-mean-field}). The sections on the fifth idea add two limitations. Their coarse states, partitions and comparison geometries are declared (\cref{asm:int-a12}), so their statements hold for the declared variables, and structure that a declared geometry contains is assumed rather than found (Supplement, \cref{rem:col-declared-geometry}). Their exact statements assume exact closure of the coarse exterior and of the coarse internal state, and outside exact closure of the exterior an approximate theorem holds only for stable linear-Gaussian collectives whose members run the consensus filter, with exact blanket conditions and a Gaussian start, and its error terms for the belief bound the upper limit and not each round; otherwise only component bounds exist (Supplement, \cref{thm:apx-h5a,open:met-approximate}). Beyond these standing conditions, several theorems are conditional on a hypothesis that their statements name: the bounds of H4 on a named contraction of the error dynamics (\cref{thm:sca-h4}), the meta-agent theorem on the hypotheses listed in \cref{sec:dis-meta} (\cref{thm:met-h5a}), the depth bound for decimated media in two and more dimensions on an inequality proved in one dimension and checked for one decimation step in two (\cref{prop:ren-depth}), and the certificate of interaction by channel interventions on its invariance conditions (Supplement, \cref{thm:idn-intervention}). The extensions of the supplement hold in the models in which they are proved, under conditions that \cref{sec:opn} states for each problem. Numerical observations hold at their stated parameters, and readings of social data, of firms and of cells are interpretations (\cref{sec:abs}; Supplement, \cref{rem:met-corporation,rem:env-interpretation}).

\subsection{Relation to the gauge-theoretic program}\label{sec:dis-gauge}

The author's gauge-theoretic manuscripts \citep{dennis2026magent,dennis2026participatory} formulate agents on associated statistical bundles over a base, with frames, connections and transport between agents as geometric structure, and the second develops a program toward physical structure. This framework takes from them only exact results that apply to linear codes in one world: the update subspace of a message and the decomposition of the meaning that a code does not transmit (\cref{prop:fw-update-subspace,prop:fw-prior-completion}); the holonomy criterion for consensus on strongly connected graphs, a frustration corollary and a tangent model of learned holonomy (\cref{thm:coh-holonomy,prop:coh-frustration}; Supplement, \cref{app:coh-consistency}); and, from the second manuscript, examples that separate small memory, sufficiency and a slow spectrum, and the lumpability form of the block-elimination identity (\cref{sec:agy-closure}; Supplement, \cref{rem:agy-minimal-memory}). The declared-selection clauses of the formation standard in the working document that accompanies the first manuscript \citep{dennis2026formation} take here the form of the null family (\cref{def:frm-target,def:agy-nulls}), and its characterization of the Gaussian location family says that mean-pooled beliefs close exactly under common-message dynamics only in that family (Supplement, \cref{rem:agy-location-family}). The extension behind \cref{hyp:int-h5} imports no result from these manuscripts. Their top-down feedback, in which, as secondary descriptions of those documents state it, a child's prior is set from its parent's belief with a blend factor and the parent is a barycenter of its children's beliefs, has as linear-Gaussian analogue literal top-down reading, whose exact alternative removes the child's own contribution from the parent's message; whether a covariance-averaging barycenter changes the size of the double count is not settled (Supplement, \cref{rem:hie-gauge}).

It does not take the geometry. Every generative model here is a model of one external physical world, and frames are optional private encodings that act only through communication (Supplement, \cref{prop:fw-frame-covariance}). True translations between encodings of one world are flat (Supplement, \cref{prop:coh-flat-true}), so holonomy measures learning or model error, not the curvature of a connection. Most results do not use frames, and a result that holds with one effective code per modality holds at identity encodings (\cref{prop:fw-effective-codes}); the results of the earlier minimal model that need frame-induced translation errors have no counterpart there and belong to the gauge-theoretic setting (Supplement, \cref{def:min-classes,rem:min-reciprocity}). In the formation model frames enter only as a fixed rotation between two groups. Whether frames that drift, which the gauge-theoretic program treats as dynamical, contribute to formation something that survives exact fusion (Supplement, \cref{conj:frm-drifting-frames}) is the question on which the two programs would meet.
